\section{The unbounded-$a$ complement of R7, including its physical zero face}
\label{sec:r7-unbounded-a}

\paragraph{Quantified source and conclusion.}
Keep the entire original weak-C2 source: $d=1$, $\lambda=\beta=0$,
the actual rational stationary coefficients, $n\ge0$, every strict gap,
the original finite stationary contact, total scale, and all prior
exclusions. In R7, $B,z,y$ are positive infinitesimals,
$a=z/B$, $t=y/B$ is a positive bounded unit, and $h,g$ are positive
infinitesimals. The original exclusion of simultaneously infinitesimal
$r$ and the source gap $j_{\rm gap}$ is retained.

\begin{theorem}
On the whole original R7 source intersection with $a$ unbounded,
the complete original parent is a finite separated source--testing
tensor sum of squares for every independent projective testing pair
and every real or complex state. Both angle signs, all literal zero
and infinite angles, literal $n=0$, and every relative opening rate
of the source infinitesimals are included. No inverse or radical of
a mixed source--testing expression is used.
\end{theorem}

The inherited bounded-$a$ R7 theorem and this theorem exhaust R7.
This statement does not admit R3, R4, or the global finite-lift theorem.

\paragraph{Actual source scale and the physical inequality.}
Use
\[
 \begin{gathered}
 s=1/a,\quad b=a^2B,\quad x=sb=z,\quad B=s^2b,\\
 w_s=ae=1+hx-g,\quad e=sw_s,\quad r=1-e,\quad l=1-h.
 \end{gathered}
 \tag{RU1}
\]
Set $K=1+x$, $R=1+hx$, $T=1+xr$ and
$j_{\rm gap}=T-l/t$. The exact original source identity is
\[
\begin{gathered}
 2r\Delta n=eQ-BK_{\rm phys},\qquad
 2r\Delta_*n/s^2=w_sQ-bK_*,\qquad \Delta_*=s\Delta,\\
 Q=1+B(1+t+ah+2ar+at)
    +B^2\{a(h+t)+a^2[h(r+1)+t(1+r-r^2)]\},\\
 K_{\rm phys}=(a+1)(h+t)(1+x),\qquad
 K_*=(1+s)(h+t)(1+x).
\end{gathered}                                                     \tag{RU2}
\]
Before bounding $b$, $Q\sim\Delta_*\sim1$ since all their nonconstant
terms contain $s$ or $x$. Here $w_s\sim1$ and $K_*\sim t$ is a
positive unit. The original inequality $eQ\ge BK_{\rm phys}$
therefore gives $b$ bounded and $n=O(s^2)$, including its literal
equality case. In particular $r\sim1$, so the original R7 exclusion
is respected. The source gap gives $t_0\ge1$; $j_{\rm gap}$ is a
unit exactly when $t_0>1$, and otherwise $t_0=1$.

The original rational formulas, after cancellation, give jointly
\[
 mg\sim\tfrac12,\quad C\sim\tfrac t2,\quad
 A_h:=Ah\sim\tfrac t2,\quad
 \beta_D:=D(1-c)=Dj_{\rm gap}/T\sim\frac1{2t},
 \quad \frac\alpha{s^3b^2}\sim\tfrac t2,
 \qquad c=\frac l{tT}.
 \tag{RU3}
\]
Thus $m,A$ are unbounded; $C$ is a positive unit. If $t_0>1$,
$D\sim1/[2(t-1)]$ is a positive unit. If $t_0=1$, $D$ is unbounded.
The canceled coefficients $mg,C,A_h,\beta_D,n\gamma,n\gamma^2$
below have unit denominators through both small-gap ends. No ratio
$h/g$, $j_{\rm gap}/h$, or comparison with a testing rate is used.

\paragraph{Four whole packets and a bounded full-parent remainder.}
Use the complete parent $P=F/4$:
\[
\begin{aligned}
 P={}&nD_uD_vX^2+mM^2+CD_v(uW-Y)^2+AD_uD_vY^2\\
 &+(D+1/2)J^2+(\alpha+1/2)D_u(vY-Z)^2
                   +D_v(X+uZ)^2/2-D_uD_vZ^2/2,\\
 M={}&vW-X+uvY-uZ,\qquad J=uW+uvX-Y-vZ.
\end{aligned}
\]
The exact coordinate inverse (RB6)--(RB7) is
\[
\begin{gathered}
 D_u=1+u^2,\quad D_v=1+v^2,\quad Z=1,\quad
 p=w_s/R,\quad q_A=D_uv\omega,\\
 N=X+u-pq_A,\qquad q_B=v(1-uX),\qquad M_B=uW-Y-cq_B,\\
 X=N+pD_uv\omega-u,\quad q_B=v(D_u-uN-puD_uv\omega),\\
 W=\omega+u(M_B+cq_B)/D_u,\qquad
 Y=u\omega-(M_B+cq_B)/D_u.
\end{gathered}                                                     \tag{RU4}
\]
The sole testing inverse is $D_u^{-1}$. Put
\[
 \begin{gathered}
 \gamma=\frac{rK}{sR},\quad \eta=Y-huW,\quad q=1-c=j_{\rm gap}/T,\\
 \mu_1=mg/R,\quad \mu_2=(mg)g/R^2,\quad n_1=n\gamma,\quad n_2=n\gamma^2.
 \end{gathered}
\]
Retain the four entire nonnegative squares
\[
 mN^2,\qquad nD_uD_v(X+\gamma v\omega)^2,
 \qquad AD_uD_v\eta^2,\qquad DM_B^2.
 \tag{RU5}
\]
All four rows vanish at the actual contact; (RB9a) verifies the first
two, and directly $Y_c=hu_cW_c$, $M_{B,c}=0$.
The $n$ packet vanishes identically on the literal $n=0$ face.

The full $A,D$ cancellations are the exact identities
\[
\begin{aligned}
 AD_uD_v(Y^2-\eta^2)&=2A_hD_uD_vuWY-A_hhD_uD_v(uW)^2,\\
 D(J^2-M_B^2)&=-2\beta_DM_Bq_B+\beta_Dq q_B^2,
 \qquad J=M_B-q q_B.
\end{aligned}                                                     \tag{RU6}
\]
Since $M=(g/R)q_A-N$, subtraction of precisely (RU5) leaves
\[
\begin{aligned}
 G={}&-2\mu_1q_AN+\mu_2q_A^2
       -2n_1D_uD_vX(v\omega)-n_2D_uD_v(v\omega)^2\\
    &+CD_v(uW-Y)^2+2A_hD_uD_vuWY-A_hhD_uD_v(uW)^2\\
    &-2\beta_DM_Bq_B+\beta_Dq q_B^2+J^2/2\\
    &+(\alpha+1/2)D_u(vY-1)^2+D_v(X+u)^2/2-D_uD_v/2.
\end{aligned}                                                     \tag{RU7}
\]
This is a whole-parent identity, including every scalar, normal and
hard coefficient. Equations (RU3), (RU6) bound all of them before
any contact restriction. In particular neither $h^{-1}$ nor
$j_{\rm gap}^{-1}$ survives.

\paragraph{One regular contact frame for the two $b$ alternatives.}
Use the pure-source parameters
\[
 \varepsilon=\frac{K\sqrt{bt}}R,\quad
 \jmath=\sqrt{s/t}\,\frac RK,\quad k=\varepsilon\jmath=\sqrt x,
 \quad E=tK/R.
 \tag{RU8}
\]
The symbol $\jmath$ is a contact scale, not $j_{\rm gap}$.
Here $\jmath$ is infinitesimal; $\varepsilon$ is a unit or
infinitesimal exactly as $b$ is. The exact inverse source chart is
\[
 x=\varepsilon^2\jmath^2,\qquad
 b=\frac{\varepsilon^2R^2}{K^2t},\qquad
 s=\frac{t\jmath^2K^2}{R^2}.
 \tag{RU9}
\]
Only source units are inverted. The contact is
$(u,\omega,v;N,M_B)=(k,\jmath,\varepsilon\jmath^2E;0,0)$.

Define
\[
\begin{gathered}
 \mathsf A=\varepsilon\omega v,\quad\mathsf B=v,
 \quad\mathsf C=Eu\omega,\quad\mathsf D=\varepsilon^2\jmath^3E,\\
 \mathsf x=\mathsf B-\mathsf C,\quad
 \mathsf y=\mathsf A-k\mathsf C,\quad
 \mathsf z=u\mathsf B-k\mathsf C,\quad
 \mathsf h=k\mathsf C-\mathsf D.
\end{gathered}                                                     \tag{RU10}
\]
The two potentially singular source combinations are
\[
 F_0=p^2/2-n_2-2n_1p+\mu_2,\qquad A_N=p-2\mu_1-2n_1.
\]
Their exact canceled quotients satisfy
\[
 F_0/b\sim t/2,\qquad A_N/(bs)\sim t.
 \tag{RU11}
\]
Both denominators are units through $s=b=h=g=0$. These assertions
follow from the original rational source, not from the inequality
$n\ge0$. In particular
\[
 F_0/\varepsilon^2=\frac{R^2}{K^2t}\frac{F_0}b,
 \qquad \alpha/\mathsf D^2=\frac1t\frac\alpha{s^3b^2},
 \qquad A_N/\varepsilon=\frac{\varepsilon R^2s}{K^2t}
                               \frac{A_N}{bs}
 \tag{RU12}
\]
are bounded without losing the zero-$n$ face.

\paragraph{Exact full scalar and both first normals.}
Put $H=\alpha+1/2$, $V_0=1-puv\omega$. At $N=M_B=0$,
$uW=\mathsf C/E+u^2c\mathsf BV_0$ and
$Y=\mathsf C/E-c\mathsf BV_0$. The complete scalar is
\[
 G_0=f_A\mathsf A^2+f_{BB}\mathsf B^2+
       f_{BC}\mathsf B\mathsf C+f_{CC}\mathsf C^2+
       \varrho\mathsf D^2,                                      \tag{RU13}
\]
with the bounded coefficients
\[
\begin{aligned}
 f_A={}&D_u^2D_vF_0/\varepsilon^2,\qquad
 \varrho=\alpha D_u/\mathsf D^2,\\
 f_{BB}={}&[Cc^2D_u^2D_v+(\beta_Dq+q^2/2)D_u^2+HD_uv^2c^2
             -A_hD_uD_v(2u^2+hu^4)c^2]V_0^2\\
          &+2HD_ucV_0-D_u/2-D_u^2\mu_2v^2\omega^2,\\
 f_{BC}={}&[2A_hD_uD_v(lu^2-1)cV_0-2HD_uv^2cV_0
                   -2HD_u+2n_1D_uD_v]/E,\\
 f_{CC}={}&[A_h(2-h)D_uD_v+HD_uv^2+D_uD_vn_2v^2]/E^2.
\end{aligned}                                                     \tag{RU14}
\]
For the full normal expansion
$G=G_0+L_NN+L_MM_B+G_{\rm hard}$, the exact coefficients are
\[
 L_N=R_A\mathsf A+R_B\mathsf B+R_C\mathsf C,
 \qquad L_M=S_B\mathsf B+S_C\mathsf C,                           \tag{RU15}
\]
where
\[
\begin{aligned}
 R_A={}&D_uA_N/\varepsilon,\\
 R_B={}&D_u(p-2n_1)v^2\omega
 -2uv[Cc^2D_uD_v+(\beta_Dq+q^2/2)D_u]V_0\\
 &+2A_hc^2uvD_vu^2(2+hu^2)V_0-2Hcuv(1+cv^2V_0),\\
 R_C={}&[2A_hcuvD_v(1-lu^2)+2Hcuv^3]/E,\\
 S_B={}&2[CcD_uD_v-(\beta_D+q/2)D_u
                      -A_hcD_vu^2(2+hu^2)]V_0
                         +2H(1+cv^2V_0),\\
 S_C={}&[2A_hD_v(lu^2-1)-2Hv^2]/E.
\end{aligned}                                                     \tag{RU16}
\]
At $u=v=0$ the complete hard matrix on $(N,M_B)$ is exactly
$\operatorname{diag}(1/2,C+1/2)$. The checker differentiates the
entire (RU7) and multiplies back (RU13)--(RU16); these are not
restricted Hessian or scalar-contact checks.

\paragraph{Bounded monic conductor, including the small-$\varepsilon$ end.}
The identities (RI14)--(RI15) apply to (RU10) without a source-shape
assumption. Their crucial inputs are checked here on the new complete
coefficients. After finite source unit denominators and $D_u^2$ are
cleared, every $\omega$ occurrence of (RU14) is through $uv\omega$
or $v^2\omega^2$, and all remaining testing powers are even in $u,v$.
The scalar raw support is exactly the five terms in (RU13).

For completeness, first replace $\mathsf A=\mathsf y+k\mathsf C$,
$\mathsf B=\mathsf x+\mathsf C$, and
$\mathsf D=k\mathsf C-\mathsf h$. The coefficient requiring a
first conductor is $S=k^2f_A+f_{BB}+f_{BC}+f_{CC}+k^2\varrho$.
Divide its contact-zero difference monically in order $u,\omega,v$.
At $u=k=\varepsilon\jmath$, the $\omega$ divided difference is in
$(\varepsilon v,v^2)$; write it as $\varepsilon U+vV$ with
$U,V\in(v)$. The last $v$ divided difference is in
$(\varepsilon,v)$ because, modulo $\varepsilon$, the remaining
$v$ dependence is even. The first $u$ divided difference restricts
to zero at $\varepsilon=u=v=0$.

These facts permit the exact bounded lift (RI15), which uses the
identity for $\mathsf Cv(\omega-\jmath)$ and never divides by
$\varepsilon$. Subtract the gradient multiple
$2k\varrho_c(-2k,1,1,-1)$. The first three contact gradients have
determinant $-\varepsilon^3\jmath^4E^2\ne0$; the CC5 adjugate
argument cancels a zero identity, and introduces no inverse of that
determinant. Each second-stage coefficient at
$\varepsilon=u=v=0$ is independent of $\omega$: the $\mathsf x$
one restricts to $2f_{BB}+f_{BC}$, and the other three to zero.
Hence every second $\omega$ divided difference is in
$(\varepsilon,v)$ and has the same bounded lift. The final repair is
\[
 k\mathsf C(\mathsf y+\mathsf z-\mathsf h-2k\mathsf x)
      =Q_1^2-\mathsf y\mathsf z.
 \tag{RU17}
\]
The pairings $Q_i$ are defined below. For each complete normal in
(RU15), one conductor suffices: after replacing $\mathsf A,\mathsf B$,
its $\mathsf C$ coefficient restricted at
$\varepsilon=u=v=0$ is independent of $\omega$, directly from
(RU16). This is precisely the normal criterion in the same proof.
When $b$ is a unit, $\varepsilon^{-1}$ is bounded and the ordinary
CC1--CC6 conductor is also available.

All contact value and first-jet hypotheses follow from the stationary
original parent and the four exact contact-zero packets (RU5).
The algorithm is a bounded, source-linear map on the stated raw
coefficient class and preserves the filtrations used next. In
particular it preserves the source ideal $(h,g)$; no separate
contact-zero assertion for its individual coefficient summands is
needed. The source substitution in each monic step uses the actual
$E$, which is a bounded unit through $h=0$.

\paragraph{The complete filtered leading model.}
Set
\[
 \Lambda=\frac1{t(1+s)},\qquad \mu=\frac1{1+s},\qquad
 T_4=\mathsf A^2+(u\mathsf B)^2+(k\mathsf C)^2+
                \mathsf D^2-4k^2\mathsf B\mathsf C.
\]
The model scalar is
\[
\begin{aligned}
 G_{\rm mod}
 &=\Lambda\mathsf x^2+T_4/2-(1/t+1/2)k^2\mathsf x^2
                      +(k^2-u^2)\mathsf C\mathsf x/t\\
 &=\left[\Lambda+\frac{u(u+k)}t-(1/t+1/2)k^2\right]\mathsf x^2
                  +T_4/2-\frac{u+k}t\mathsf x\mathsf z.
\end{aligned}                                                     \tag{RU18}
\]
The second equality is the exact conductor identity
$\mathsf C(u-k)=\mathsf z-u\mathsf x$. Use the first expression
for the raw coefficient filtration, since it retains the even-$u$
class. Its angular raw coefficients are
\[
\begin{aligned}
 F_{BB}&=\Lambda+u^2/2-(1/t+1/2)k^2,\\
 F_{BC}&=-2\Lambda-u^2/t+(3/t-1)k^2,\\
 F_{CC}&=\Lambda+u^2/t-2k^2/t.
\end{aligned}                                                     \tag{RU19}
\]
At $h=g=0$, exact source collection gives the following two
coefficientwise filtrations:
\[
\begin{array}{c|c|c}
 &b\text{ a unit}&b\text{ infinitesimal}\\ \hline
 \text{variable weights}&(k,u,\omega,v)=(1,1,1,2)
     &(\varepsilon,\jmath,u,\omega,v)=(1,1,2,1,3)\\
 \operatorname{wt}(\mathsf x)&2&3\\
 \operatorname{wt}(\mathsf y,\mathsf z,\mathsf h,Q_1,Q_2)&3&5\\
 f_A-1/2,\ \varrho-1/2&\ge2&\ge2\\
 f_{BB}-F_{BB},\ f_{BC}-F_{BC},\ f_{CC}-F_{CC}&\ge4&\ge6\\
 S_B-\mu,\ S_C+\mu&\ge2&\ge4\\
 (R_A,R_B,R_C)&\ge(2,3,3)&\ge(3,5,5).
\end{array}                                                       \tag{RU20}
\]
These statements concern the chosen raw coefficients separately;
polynomial cancellation between raw products is not used to infer
them. Consequently
\[
\begin{array}{c|c|c}
 &b\text{ a unit}&b\text{ infinitesimal}\\ \hline
 G_0-G_{\rm mod}&O_{\rm wt}(8)&O_{\rm wt}(12)\\
 L_M-\mu\mathsf x&O_{\rm wt}(4)&O_{\rm wt}(7)\\
 L_N&O_{\rm wt}(5)&O_{\rm wt}(8).
\end{array}                                                       \tag{RU21}
\]
Here each equality is modulo the source ideal $(h,g)$.

To reproduce the collection, the necessary $b$-infinitesimal source
jets at $h=g=0$ are
\[
\begin{aligned}
 C&=t/2-(t+1)s/2+(t+1)s^2/2+O_{\deg(s,b)}(3),\\
 A_h&=t(1-s+s^2)/2+O_{\deg(s,b)}(3),\\
 \beta_D&=[1-s+s^2-(t+1)sb]/(2t)+O_{\deg(s,b)}(3),\\
 n_1&=(s-s^2-tsb)/2+O_{\deg(s,b)}(3).
\end{aligned}                                                     \tag{RU22}
\]
Also $c=(1-k^2)/t+O_{\rm wt}(6)$,
$E=t(1+k^2)$ at $h=0$, and
$s=t\jmath^2+O_{\rm wt}(6)$. For the unit-$b$ filtration use
$s=k^2/b$ and the respective coefficients of $s$:
$-(t+1)/2,-t/2,-[1+(t+1)b]/(2t),(1-bt)/2$.
The source cancellations (RU11)--(RU12) supply the first and last
rows of (RU20). The checker verifies all needed source derivatives,
each angular raw coefficient through weight two or four, the whole
scalar through weight six or ten, and each complete $M_B$ normal
through weight three or five. Scalar weights are even in both
charts. The $M_B$ weights are even in the first and odd in the
second, giving exactly the next weights in (RU21).

Apply the bounded conductor to the exact differences in (RU21).
Their coefficientwise raw filtration makes every correction Gram
coefficient have strictly positive weight. Terms in $(h,g)$ lift
boundedly and remain infinitesimal at arbitrary relative rates.
The whole hard block differs infinitesimally from
$\operatorname{diag}(1/2,(t+1)/2)$, since all its coefficients are
bounded and $u,v,h,g,s$ are infinitesimal.

On the negative-core contact chart take $u,v,\omega>0$. The separated
nonnegative monomials
\[
 a_i^2=(\mathsf A,u\mathsf B,k\mathsf C,\mathsf D),\qquad
 Q_1=a_1a_2-a_3a_4,\quad Q_2=a_1a_3-a_2a_4
\]
have the exact product relation yielding
\[
 T_4=\tfrac12\{(\mathsf y-\mathsf z)^2+\mathsf h^2+
                 \mathsf y^2+(\mathsf z+\mathsf h)^2\}
                  +Q_1^2+Q_2^2.                                \tag{RU23}
\]
Its three-dimensional difference matrix has determinant $1/2>0$.
Each radical separates into source and testing factors.
The full eight-root residue on
$(\mathsf x,\mathsf y,\mathsf z,\mathsf h,Q_1,Q_2,N,M_B)$
therefore consists of half this strict circuit, the $N$ pivot $1/2$,
and the complete coupled block
\[
 \begin{pmatrix}1/t&1/2\\1/2&(t+1)/2\end{pmatrix},\qquad
 \det=\frac{t+2}{4t}>0.                          \tag{RU24}
\]
The additional cross term in the second line of (RU18) has
infinitesimal coefficient. WA10 now gives a finite separated SOS
for the exact remainder. Restore all four actual squares (RU5).
This proves the whole local contact chart in both $b$ alternatives.

\paragraph{The $\omega$-unit state chart.}
If $u,v$ are infinitesimal, $Z=1$, and $\omega$ is a positive unit,
replace the retained centered $n$ packet by $nD_uD_vX^2$ itself.
The complete remainder is (RU7) with the two $n_1,n_2$ terms removed.
Use the frame $(A_o,\mathsf x,\mathsf D,N,M_B)$, $A_o=\omega v$.
The exact identities
$\mathsf A=\varepsilon A_o$, $\mathsf B=A_o/\omega$,
$\mathsf C=A_o/\omega-\mathsf x$ use only a pure-testing unit
inverse. In particular the first scalar coefficient becomes
$D_u^2D_v(p^2/2+\mu_2)A_o^2$, and the first $N$ coefficient
becomes $D_u(p-2\mu_1)A_o$; the printed $\varepsilon$ denominators
cancel completely. The full residue has the $A_o,\mathsf D,N$
pivots $1/2,1/2,1/2$ and the same block (RU24).
The complete normal residues are $0,\mathsf x$. Hence this chart
also has a strict finite separated Gram.

\paragraph{A remote cap lemma and complete bounded parents.}
Away from the contact chart there is no need for centered excesses.
Retain the original nonnegative $n,\alpha$ packets and the positive
excesses of $m,A$ above chosen positive real caps $m_c,A_c$.
If $D$ is unbounded, retain its original excess above a real cap
$D_c$ too; otherwise leave $D$ at its actual bounded value.
The remainder is the full original parent with these three
coefficients capped, and all its coefficients are bounded.

We use the elementary finite-dimensional fact that a real quadratic
form positive definite on the common kernel of finitely many real
linear rows becomes positive definite after sufficiently large
positive real multiples of the squared rows are added. Choose
coordinates on the common kernel and its orthogonal complement;
the assertion follows by a Schur complement. Unbounded actual
$m,A,D$, when used, exceed each chosen real cap. Thus every cap
is an exact nonnegative packet subtraction in the original parent,
and a strict residue gives WA10 for the entire bounded remainder.
The real caps depend only on real testing residues, never on a
source/testing rate comparison.

\paragraph{Finite remote $v$ with nonzero residue.}
Write $U=u^2,V=v^2$ for real testing residues, with $V>0$.
On the common kernel $Y=M=0$, $X=vW-uZ$, and the remaining
two-state form is exactly
\[
 CD_vU W^2+(D+1/2)(uD_vW-vD_uZ)^2
                  +VD_vW^2/2-D_uVZ^2/2.                         \tag{RU25}
\]
Its determinant is
\[
 VD_uD_v\{DV/2+U[CD-(D+1/2)/2]+C(D+1/2)U^2\}.
 \tag{RU26}
\]
For $t_0>1$, insert $C=t/2,D=1/[2(t-1)]$; it equals
$D_uD_v[V^2+t^2U^2V]/[4(t-1)]>0$.
Its $W^2$ coefficient is positive since $V>0$.
For $t_0=1$, $C=1/2$ and one may choose a real $D$ cap: the braces
are $D(V+U^2)/2+(U^2-U)/4$, positive for sufficiently large $D$.
The cap lemma then chooses $m_c,A_c$ for a strict full state matrix.

\paragraph{Finite unit $u$ with infinitesimal $v$, including zero.}
Use $N=-M=X+uZ-v(W+uY)$, $R_t=uW-Y$, and $V_Z=vZ$.
After the original $n,\alpha$ packets and allowed positive excesses
are retained, the entire remainder lies in the square of the four
roots $(N,R_t,Y,V_Z)$. Only the pure-testing unit $u$ is inverted
to write $W=(R_t+Y)/u$; no inverse of $v$ is needed. Its exact
residue is
\[
\begin{aligned}
 (m_c+1/2)N^2+CR_t^2+A_cD_uY^2
 +(D+1/2)(R_t-D_uV_Z)^2-D_uYV_Z-D_uV_Z^2/2.
\end{aligned}                                                     \tag{RU27}
\]
For $t_0>1$, its restriction to $Y=0$ has determinant
$t^2D_u(D_u-1)/[4(t-1)]>0$ on $(R_t,V_Z)$.
Choose $A_c$ sufficiently large. For $t_0=1$, the common kernel
$Y=0,R_t-D_uV_Z=0$ has value $D_uu^2V_Z^2/2>0$;
large real $D_c,A_c$ give strictness. The normal pivot is positive.
The full four-root support and the residue are checked directly;
literal $v=0$ is their polynomial value.

\paragraph{Both finite angles infinitesimal, and all state charts.}
The exact SD5 pure-core sign identity deals with the nonnegative-core
branch and literal zero axes by original-event squares. On the
negative branch use its standard gauges $u,v>0$ and normalize the
state by a maximal coordinate. If $Z$ is infinitesimal, $W$ is a
unit and the exact BE4 selected four-root identity applies to
$(X,Y,uW,vW)$, with positive caps of $m,A$ provided by (RU3).
All other original events remain as nonnegative packets.
If $Z$ is a unit normalize $Z=1$; the same sign identity gives
bounded positive $\omega$. Its infinitesimal and unit alternatives
are exactly (RU18)--(RU24) and the subsequent $\omega$-unit chart.

\paragraph{Every reciprocal chart and literal projective infinity.}
For unbounded $u$, put $p_u=1/u$ and use the entire homogeneous
polynomial $p_u^2P$. If bounded $v$ has nonzero residue, the full
capped four-state residue is
\[
 CD_vW^2+A_cD_vY^2+(D_c+1/2)(W+vX)^2
                          +(m_c+1/2)(vY-Z)^2\succ0.             \tag{RU28}
\]
Here $D_c$ denotes either a positive real cap or the actual positive
unit residue. The four square rows have common kernel zero because
$v$ is a unit. If instead $v$ is infinitesimal, including zero,
use the exact BE $(\infty,0)$ selected full frame
$(p_uX+Z,W,Y,vX)$ with positive caps of $C,A$. Its strict residue
is the sum of positive multiples of
$(p_uX+Z)^2,(W+vX)^2,W^2,Y^2$; the other original packets are
retained. This handles the zero-$n$ degeneration at that corner.

For unbounded $v$ and bounded $u$, put $q_v=1/v$. The full capped
residue of $q_v^2P$ is
\[
 m_c(W+uY)^2+C(uW-Y)^2+(A_c+1/2)D_uY^2
                         +D_c(uX-Z)^2+D_uX^2/2\succ0.           \tag{RU29}
\]
Its common square kernel is zero for every bounded real $u$,
including zero. Finally, if both angles are unbounded, use
$p_u^2q_v^2P$ and the full frame
$(X,W,Y,P_Z=p_uZ,Q_Z=q_vZ)$. The exact residue is
\[
 CW^2+A_cY^2+D_c(X-P_Z)^2+X^2/2
                           +(m_c+1/2)(Y-Q_Z)^2\succ0.           \tag{RU30}
\]
The two nontrivial block determinants are $D_c/2$ and
$A_c(m_c+1/2)$, both positive. The entire polynomial has this
five-root support, with no inverse of $p_u$ or $q_v$ after
homogenization. All three formulas are checked from the original
parent. Their literal reciprocal zeros are polynomial values;
restoring a nonzero homogeneous factor uses a pure-testing square.

The charts exhaust independent testing pairs and states. Apply each
real identity to real and imaginary parts of a complex state, restore
every retained actual packet and separated clearing factor, and use
the exact sign gauges for both angle signs. This completes the stated
all-tests theorem without a restriction on any opening ratio.

\paragraph{Dependencies and the scope of the checks.}
The exact dependency manifest is
\texttt{dependencies\_r7\_unbounded\_a.json}. It distinguishes the
actual source definitions, full-parent identity, algebraic conductor,
sign/state reduction, and strict finite Gram lemma from the source
theorems surrounding those identities. The companion
\texttt{check\_r7\_unbounded\_a.py} reconstructs source and whole
parent expressions, checks complete scalar and normal allocations,
the two coefficientwise filtrations, and the strict local and remote
matrices. The separate-review checker
\texttt{review\_check\_r7\_remote.py} also reconstructs the complete
four-root and five-root supports without dividing by a zero testing
coordinate. These are exact symbolic corroborations; the monic
coefficient-class argument and exhaustive joins above supply the
mathematical proof. No formal verification or external peer review
is claimed.
